\section{Design}
\label{sec:design}

\sfs{} is organized around a logical unit rather than a particular API.  A unit
may be presented as a file, a metadata-only directory, an embedded blob, or an
application object, but it has the same identity, version state, encryption,
and publication rules in every case.  This section first defines that state and
then follows it across local storage, synchronization, and the encrypted service
boundary.

\subsection{Units, Streams, and Fragment-Causal State}
\label{subsec:mvcc}

A unit has a stable 128-bit UUID and up to two independently versioned streams:
content and filesystem metadata.  The metadata stream contains attributes such
as mode, ownership, timestamps, symlink target, flags, and extended attributes.
Separating the streams prevents a concurrent metadata-only operation from
becoming a content conflict.  Paths are mutable keys, not identities.  A key
catalog maps raw path bytes to UUIDs, while an identity catalog maps UUIDs to
the current unit-record location.  Rename changes the former.  Relocation
changes the latter.  Several keys may therefore name one UUID, providing the
storage primitive used for aliases and hard links.

For each present stream, the current unit record contains an ordered
\emph{unit map} $M$, a location vector $L$, a sparse version vector $VV$, the
fragment-size exponent $e$, and the logical length of the final fragment.  The
entry $M[f]$ is a packed causal dot $(a,c)$: replica alias $a$ and that replica's
monotone counter $c$.  The corresponding $L[f]$ locates the selected physical
encoding.  A write increments the local component of $VV$ and assigns the new
dot only to touched fragments.  Untouched positions retain their dots and
locations.

Each new unit record points to its predecessor.  Resolving an old logical
version walks this parent chain and chooses, for each position, the newest
fragment version no later than the requested version.  The current read path
does not perform this walk: the head record already contains the complete map
and locations of the materialized current unit.  Thus history is native but is
kept off the common read path.

Named commits are cuts over selected units.  Per-unit bitmaps pin the fragments
needed by a cut so retention does not reclaim them.  Unnamed history may be
thinned by the configured retention policy.  The guarantee is therefore
retention-scoped rather than an assertion that every superseded byte remains
forever.  Root-key rotation introduces an additional epoch boundary discussed in
\cref{sec:discussion}.

\subsection{Deterministic Fragment Geometry}
\label{subsec:chunking}

Content-defined chunking preserves boundaries after insertions and is effective
for deduplication, but it makes boundaries content-dependent and requires a
scan to rediscover them.  \sfs{} instead uses fixed positions within a unit.
For fragment size $2^e$, byte offset $x$ maps to fragment

\begin{equation}
  f(x) = x \mathbin{\mathtt{>>}} e,
\end{equation}

which is constant time and identical in the Rust and C implementations.  The
production schedule is a monotone exponent staircase:

\begin{table}[t]
  \centering
  \small
  \caption{Current v12 fragment-size schedule.  The last band is capped.  Its
    fragment count consequently grows linearly for sufficiently large units.}
  \label{tab:geometry}
  \begin{tabular}{@{}lll@{}}
    \toprule
    Logical unit size & Fragment size & Fragment count near band top \\
    \midrule
    $<16$ KiB & 4 KiB & 1--4 \\
    16--256 KiB & 16 KiB & 1--16 \\
    256 KiB--64 MiB & 256 KiB & 1--256 \\
    $\geq64$ MiB & 4 MiB & $\geq16$ \\
    \bottomrule
  \end{tabular}
\end{table}

The staircase approximates square-root growth before reaching the 4-MiB cap.
It reduces per-fragment metadata for large units while retaining small update
granularity for small files.  A unit crossing a band boundary is re-fragmented.
That transition loses delta reuse across the old and new geometry: unpinned old
fragments are freed, whereas commit-pinned fragments are retained as history.
Manifest construction and comparison remain linear in the number of units and
fragments.  Only offset-to-fragment lookup is $O(1)$.  In steady state, payload
transfer is proportional to changed or missing fragments, not to total retained
history.

\subsection{Atomic Publication and Recovery}
\label{subsec:commit}

The first two 4-KiB blocks of a container are alternating header slots.  Each
v12 header contains the format and cipher configuration, catalog roots,
publication sequence, writer and key epochs, eviction-tail watermark, and
password-KDF salt.  CRC32 detects torn or random corruption.  HMAC-SHA256 under a
root-key-derived header key authenticates security-relevant fields.  Of the
supported, CRC- and MAC-valid slots, the one with the highest commit sequence is
active.

A normal mutation writes fresh content fragments, a new unit record, and
copy-on-write catalog spines.  It then flushes these blocks and writes the
inactive header slot with the new roots and a higher sequence.  Before this
publication, all new objects are unreachable from the active roots.  After it,
the new state is complete.  A power failure therefore selects either the old
or the new root set on reopen, not a mixture.  Several operations, including
record-plus-index updates in the annex, may be staged under one publication.

Atomic root publication is not a claim that the system has no logging.  The
optional asynchronous mount path uses a WAL, and bounded undo records protect
the deliberately limited in-place overwrite paths.  Recovery replays only
records newer than the header's applied sequence.  If catalog roots themselves
are unavailable, an offline scanner recognizes self-describing unit and
eviction records, reconstructs identity heads, preserves recoverable paths, and
places otherwise unnamed units under a lost-and-found keyspace.  When several
orphan candidates have no parent relation, highest physical address is only a
recovery heuristic, not a causal proof of recency.

\subsection{Three-Region Allocation and Retention}
\label{subsec:alloc}

The container uses a static catalog-front layout: catalog blocks grow forward,
current unit data occupies the middle, and retained/evicted history grows
backward from the tail.  Current file content is not stored below one global
copy-on-write tree, so changing a fragment does not recursively shadow a
content-tree path to a filesystem root.  The key and identity catalogs are
copy-on-write tries and still pay their own spine-update cost.  Transactions
therefore open a reclaim scope that recycles intermediate, as-yet-unpublished
catalog spines during a batch.

\begin{figure}[t]
  \centering
  \begin{tikzpicture}[x=1cm,y=0.85cm,font=\small]
    \draw[fill=black!8] (0,0) rectangle (0.65,1);
    \draw[fill=black!8] (0.65,0) rectangle (1.3,1);
    \draw[fill=RoyalBlue!12] (1.3,0) rectangle (3.0,1);
    \draw[fill=ForestGreen!12] (3.0,0) rectangle (5.0,1);
    \draw[fill=white] (5.0,0) rectangle (6.3,1);
    \draw[fill=BurntOrange!15] (6.3,0) rectangle (8.3,1);
    \draw[fill=black!5] (8.3,0) rectangle (9.5,1);

    \node[font=\scriptsize] at (0.325,0.5) {H0};
    \node[font=\scriptsize] at (0.975,0.5) {H1};
    \node[align=center] at (2.15,0.5) {catalog\\head};
    \node[align=center] at (4.0,0.5) {live unit\\data};
    \node[font=\scriptsize] at (5.65,0.5) {free};
    \node[align=center] at (7.3,0.5) {retained\\history};
    \node[font=\scriptsize] at (8.9,0.5) {WAL};

    \draw[->,>=Latex,thick] (1.4,1.35) -- (4.9,1.35)
      node[midway,above] {forward allocation};
    \draw[->,>=Latex,thick] (8.2,1.35) -- (6.4,1.35)
      node[midway,above] {backward growth};
    \node[anchor=east,font=\footnotesize] at (0,0.5) {byte 0};
    \node[anchor=west,font=\footnotesize] at (9.5,0.5) {container end};
  \end{tikzpicture}
  \caption{Schematic v12 physical layout, not to scale.  Two fixed 4-KiB
  header slots alternate as the publication point.  Catalog and live regions
  allocate from the low end, while retained history grows down from the
  high end.  Their gap is available capacity.  When asynchronous writes are
  enabled, a fixed WAL reservation occupies the top of the container and the
  history frontier grows below it.}
  \label{fig:container-layout}
\end{figure}

\Cref{fig:container-layout} makes the independent growth directions and
publication slots explicit.  A collision between the forward and backward
frontiers grows the container rather than allowing the regions to overlap.

Superseded current fragments become self-describing history records when they
must survive.  Retention can thin or explicitly evict unpinned history, and
released ranges can be trimmed.  Limited defragmentation relocates only state
whose history and pin relationships remain valid.  This is not a wandering
log and has no background segment cleaner, but it does have explicit retention,
reclamation, defragmentation, and trim operations.

\subsection{Version-Vector Synchronization and Strains}
\label{subsec:sync}

Synchronization first compares unit version vectors and fragment-dot maps.  For
one fragment position $f$, let $d_L=M_L[f]$ and $d_P=M_P[f]$ be the local and
peer dots.  \sfs{} classifies the position as follows:

\begin{equation}
\operatorname{choose}(f)=
\begin{cases}
L, & d_L=d_P,\\
L, & VV_L\text{ has observed }d_P,\\
P, & VV_P\text{ has observed }d_L,\\
\mathit{conflict}, & \text{otherwise.}
\end{cases}
\end{equation}

The first case is identical state.  In the next two cases, one version vector
witnesses the other dot, so choosing the witnessing side selects a causal
descendant rather than discarding concurrent information.  In the final case
neither side witnesses the other dot.  Selecting either would lose an unseen
update, so the safe result is conflict retention.  Causal dominance and version
vector join are independent of which replica is called local, idempotent, and
commutative.  Replicas presented with the same frontier therefore select the
same non-conflicting fragments and joined vector.

If no position is classified as a conflict, the result selects the indicated
fragments and joins the version vectors.  In particular, two replicas that
modify disjoint fragments from a shared base converge without a
content-specific merge algorithm.  If any position contains mutually unseen
dots, \sfs{} preserves both unit heads as concurrent \emph{strains}.  This is a
same-fragment write conflict, not necessarily a byte collision: two disjoint
byte ranges inside one fragment still produce a strain.  Mounted clients expose
the condition through the \texttt{user.sfs.conflict} extended attribute and the
tooling provides explicit resolution operations.

\Cref{fig:fragment-merge} illustrates the disjoint case and its same-position
boundary.

\begin{figure}[t]
  \centering
  \begin{tikzpicture}[
    frag/.style={draw, rounded corners=1pt, minimum width=1.3cm,
                 minimum height=0.55cm, font=\footnotesize},
    changeda/.style={frag, fill=RoyalBlue!18},
    changedb/.style={frag, fill=ForestGreen!18},
    group/.style={draw, dashed, rounded corners, inner sep=3pt},
    flow/.style={->, >=Latex, thick}]
    \node[frag] (base0) at (0.7,2.6) {$f_0:d_0$};
    \node[frag] (base1) at (2.15,2.6) {$f_1:d_1$};
    \node[frag] (base2) at (3.6,2.6) {$f_2:d_2$};
    \node[group, fit=(base0)(base2), label=above:{shared base}] (base) {};

    \node[changeda] (a0) at (-1.85,0.9) {$f_0:a_3$};
    \node[frag] (a1) at (-0.4,0.9) {$f_1:d_1$};
    \node[frag] (a2) at (1.05,0.9) {$f_2:d_2$};
    \node[group, fit=(a0)(a2), label=left:{replica A}] (agroup) {};

    \node[frag] (b0) at (3.25,0.9) {$f_0:d_0$};
    \node[changedb] (b1) at (4.7,0.9) {$f_1:b_4$};
    \node[frag] (b2) at (6.15,0.9) {$f_2:d_2$};
    \node[group, fit=(b0)(b2), label=right:{replica B}] (bgroup) {};

    \node[changeda] (m0) at (0.7,-0.9) {$f_0:a_3$};
    \node[changedb] (m1) at (2.15,-0.9) {$f_1:b_4$};
    \node[frag] (m2) at (3.6,-0.9) {$f_2:d_2$};
    \node[group, fit=(m0)(m2), label=below:{merged head}] (merged) {};

    \draw[flow] (base.south) -- (agroup.north);
    \draw[flow] (base.south) -- (bgroup.north);
    \draw[flow] (agroup.south) -- (m0.north);
    \draw[flow] (bgroup.south) -- (m1.north);
  \end{tikzpicture}
  \caption{Fragment-causal merge from a shared base.  A replaces position
  \(f_0\) and B independently replaces \(f_1\).  At each changed position the
  writer's new dot causally dominates the other replica's base dot, so the
  merged head contains both updates.  If both replicas instead placed mutually
  unseen dots at the same position, the final case of
  \(\operatorname{choose}\) would retain the two unit heads as strains.}
  \label{fig:fragment-merge}
\end{figure}

We use \emph{strain} for a coexisting head of one logical unit.  Unlike a
repository branch, it does not name a repository-wide line of development.
Unlike a server fork, it does not imply equivocation by the transport.  It is
the retained, user-resolvable consequence of mutually unseen dots at one or
more positions.

Record projections are transferred before or alongside their missing payload
fragments.  A receiving engine may temporarily represent a remote-selected
fragment as a hole and fill it from the transport.  The transport retains a
frontier of pairwise-concurrent record projections rather than a single
last-writer-wins value, so a third replica can discover every unresolved side
from any peer that has observed it.

\subsection{One Protocol for SaaS and Live Peers}
\label{subsec:transport}

All sync intelligence remains in the engine behind a transport interface.  The
interface lists unit frontiers, tests and transfers fragment objects keyed by
$(\mathit{account},\mathit{uuid},f,d_f)$, exchanges encrypted record
projections, and carries capability, writer-set, and key-grant objects.  It does
not expose filesystem operations to the service.

The hosted implementation is a multi-tenant opaque store.  It authenticates
accounts with SRP-6a, derives account scope from the bearer token rather than a
client-supplied path, and persists opaque blocks and record frontiers.  Version
vectors remain visible because the service maintains causal frontiers.  The
service also observes account identifiers, unit UUIDs, fragment indices and
versions, object sizes and counts, timing, and access patterns.  It does not
receive VFS paths, file metadata plaintext, content plaintext, or the client
root key in an encrypted profile.

A live peer implements the same interface directly over an open \sfs{} engine.
Transport calls become manifest projection, fragment export/import, and record
frontier export/import.  No second reconciliation algorithm is introduced.  A
network peer authenticates possession of a root-key-derived P2P key using a
single-use challenge.  Key grants are themselves container units and therefore
propagate through hosted or peer topologies.  The service is thus a replaceable
availability and rendezvous component rather than the authority for filesystem
semantics.

\subsection{Encryption and Multi-Writer Authorization}
\label{subsec:crypto}
\label{subsec:multiwriter}

A password is stretched with Argon2id into a 32-byte root key.  Domain-separated
derivations produce the metadata key, header-MAC key, and one container-level
key for each cryptographic content suite: 32 bytes for GCM and 64 bytes for XTS
(two AES-256 halves).  These keys are suite-specific but not fragment-specific.
Fragment context supplies the nonce or tweak.  Unit records, catalog nodes,
filesystem metadata, paths, writer sets, and key material stored in the
container are AES-GCM protected in v12.  Content is
encoded per fragment as AES-256-GCM, AES-256-XTS, or plaintext.  GCM provides
confidentiality and authentication.  XTS is length-preserving confidentiality
only.  Plaintext is an explicit benchmark/trusted-media mode.  Fragment context
binds UUID, position, causal version, and key epoch into nonce/tweak derivation.

\begin{table}[t]
  \centering
  \small
  \caption{Content-suite claim boundary in v12.  Metadata remains AES-GCM in
  every row.  WriterSet is an additional requirement for the secure hosted
  profile.}
  \label{tab:cipher-profiles}
  \begin{tabular}{@{}llll@{}}
    \toprule
    Content & Confidential & Authenticated & Intended role \\
    \midrule
    GCM & yes & yes & secure hosted/application profile \\
    XTS & yes & no & compatibility and performance mode \\
    \textsc{none} & no & no & trusted-media and benchmark mode \\
    \bottomrule
  \end{tabular}
\end{table}

The paper's end-to-end-encrypted deployment profile means GCM content plus
WriterSet signing.  In that profile neither the hosted service nor its at-rest
database receives the client root key or decryptable VFS state.  TLS protects
the wire hop, but confidentiality from the service follows from the client-side
container encryption, not from TLS.  This is client-side-encrypted storage
rather than cryptographic zero knowledge because the synchronization metadata
listed above remains visible.

WriterSet containers have an Ed25519 owner key and an owner-signed, monotone
writer-set object.  Members sign the replica-invariant logical portion of a unit
record: UUID, fragment dots, version vectors, geometry, and an optional
application head.  Replica-local locations, parent pointers, pins, cipher
encodings, and strain addresses are excluded so relocation, re-ciphering, and
cross-replica import do not invalidate signatures.  The signature consequently
authorizes and attributes a logical record frontier.  It is not a per-fragment
provenance proof or a content digest.  GCM authenticates content in the secure
profile.

Removal is represented by a tombstone and is coupled to a higher key epoch.
Root-key rotation re-encrypts live content, metadata, catalogs, and unresolved
strains before atomically publishing the new epoch.  New root keys are sealed
to retained members with ephemeral X25519/HKDF/AES-GCM grants.  A peer that has
not yet received its matching grant remains safely on its old epoch.  A retained
peer can adopt the new key and writer set in one local publication.  Rotation is
forward-only in v12: pre-rotation parent history and retention pins are not
carried across the key boundary, while live conflict strains are.
